\section*{The Conflict[s] in Nigeria}
\textcolor{red}{I think we want something, either at the beginning or end of this section that summarizes what we want people to take away from this account: lots of violent groups, fight eachother and government, civilian victimization plays roll in driving conflict.}

The year 2000 marks the beginning of a critically tense transition period in Nigeria. Following the 1999 presidential election of Olusegun Obasanjo, Nigeria's 15-year long military dictatorship ends. Meanwhile, a resurgent Islamic political movement consolidates influence in Nigeria's predominantly Muslim northern states. Sharia Law, a penal code in effect for hundreds of years in Nigeria but disbanded in 1960, is re-implemented following the 2000 elections. The implementation of Sharia law then leads to numerous riots and clashes, largely between Muslims and Christians (one clash, in Nigeria's second largest city of Kano, ends with over 100 dead in October of 2001).\footnote{See \url{http://articles.latimes.com/2001/oct/16/news/mn-57819}.}

President Obasanjo faces many more violent episodes during his first term. As Christian-Muslim disputes continue in the Northern regions, a tribal war foments in the eastern-central state of Benue. Thousands of people are forced to flee the area, and troops reportedly target unarmed civilians in retaliation for the abduction of nearly 20 soldiers. Obasanjo's second term (beginning in 2003) is marked by continuing Christian-Muslim clashes in the north,  and deepening violence in the Niger Delta where militants vie control of oil production centers. 

In 2002, Mohammed Yusuf emerged as the radical leader of the religious extremist group, the Boko Haram.\footnote{There is evidence that the early leadership of the movement formed at this time, though in many reports the fully armed, organized, radical group forms closer to the year 2008. See \citet{walker:2016} for details.} Yusuf spreads an anti-state, anti-elite ideology couched in Sunni Muslim teachings. Conflict in the northern region of Nigeria continues along religious lines, and by 2008 Boko Haram has grown increasingly well organized, effectively operating as a mini-state with institutions, a welfare system, and religious police.  Clashes between Boko Haram and state security forces gain momentum as the 2009 uprising begins in the state of Bauchi and spreads to Borno, Yobe, and Kano. These battles are the first sustained conflict between Boko Haram and the government, and culminates in the death of scores of police officers, more than 700 members of Boko Haram, and the capture of Mohammed Yusuf.\footnote{See \url{http://edition.cnn.com/2014/06/09/world/boko-haram-fast-facts/}.}

In 2010, Vice President Goodluck Jonathan succeeds to the presidency after the death of his predecessor. The death of Mohammed Yusuf, killed by police in a raid, leads to an intensification of tactics by Boko Haram, including greater violence against civilians. Demonstratively, in December of 2010, Boko Haram bombs kill an estimated 80 people in the central city of Jos. This event sparks a response from opposing Christians leading to roughly 200 more deaths.\footnote{See \url{http://www.aljazeera.com/news/africa/2013/09/201397155225146644.html}.}

Attacks on civilians continue and in April of 2014 Boko Haram chooses a target that sparks international condemnation: 276 Chibok schoolgirls. Though this is not the first time that civilians organize against the insurgency's violent tactics, high levels of protests and vigilantism ensue following the children's kidnapping. %This devastating attack even sparks the now internationally known ``\#bringbackourgirls'' social media hashtag \citep{maiangwa:agbiboa:2014}. 
Violence against civilians continues to displace and disappear civilians.  In March 2014, the Director-General of the Nigerian National Emergency Management Agency (NEMA) reported that more than 250, 000 people were displaced as a result of the fighting in north-eastern Nigeria. %In the same month, the Boko Haram detonate an improvised explosive devise (IED) in a market in Ngurosoye village, killing at least 16 people. 
% \textcolor{red}{add AI report cite.}

Battles between Boko Haram, local militant groups, and security forces continue today. While Boko Haram is the dominant armed non-state actor, there are also major ethnic militias such as the Urhobo Ethnic Militia and the MEND.\footnote{MEND is one of the largest militant groups in the Niger Delta region. The militant organization is expressly concerned with the public and private production of oil in the region.} Yet, Boko Haram's entrance into the conflict network in 2009 has had unique consequences. 

Figure~\ref{fig:nigMap} shows the spatial distribution of violence in Nigeria from 2001 to 2016, points in red represent conflictual events in which Boko Haram is involved (either as the initiator or target) while points in blue designate events that do not involve Boko Haram. What becomes clear is that Boko Haram's entrance into the system corresponds with an increase in both the level and spread of conflict. %Interestingly, this proliferation of conflict is not attributable to just armed conflict events involving Boko Haram, we see that the likelihood of conflictual events between actors on the other side of the country substantially increase.

\begin{figure}[ht]
	\includegraphics[width=.95\textwidth]{nigConfMap.pdf}
	\caption{Spatial distribution of conflict in Nigeria from 2001 to 2016. Points colored in blue represent conflict events between sets of actors not including Boko Haram, and points in red events that do involve Boko Haram as either the initiator or target of conflict. }
	\label{fig:nigMap}
\end{figure}
\FloatBarrier


